\begin{examplebox}
	\textbf{\textcolor{CExample}{例1（基础）}}
	\textbf{\textcolor{CExample}{题目：}}
	已知圆柱底面半径 $R=3$，高 $h=6$，判断该圆柱是否存在内切球，若存在，求内切球半径。
	\textbf{\textcolor{CExample}{解题步骤：}}
	\begin{description}[style=nextline, leftmargin=2.8em, labelsep=0.5em, itemsep=0.45em]
		\item[\textbf{第一步（验证条件）：}]
			验证高与底面直径的关系。
			\[
				h = 2R.
			\]
			\par\textbf{理由：} 代入 $R=3$，得 $2R=6$，等于高 $h=6$，故 $h=2R$ 成立。
		\item[\textbf{第二步（求半径）：}]
			由定理，内切球半径等于底面半径。
			\[
				r = 3.
			\]
			\par\textbf{理由：} 因为 $r=R$，且 $R=3$，所以 $r=3$。
		\item[\textbf{第三步（作答）：}]
			该圆柱存在内切球，半径为 $3$。
			\[
				r = 3.
			\]
			\par\textbf{理由：} 结论。
	\end{description}
	\textbf{\textcolor{CExample}{关键结论：}}
	\textbf{若 $h=2R$，则存在内切球且 $r=R$。}

	\medskip
	\textbf{\textcolor{CExample}{例2（稍变形）}}
	\textbf{\textcolor{CExample}{题目：}}
	已知圆柱存在内切球，且内切球体积为 $\frac{32}{3}\pi$，求圆柱的底面半径和高。
	\textbf{\textcolor{CExample}{解题步骤：}}
	\begin{description}[style=nextline, leftmargin=2.8em, labelsep=0.5em, itemsep=0.45em]
		\item[\textbf{第一步（列方程）：}]
			设球半径为 $r$，由体积公式列出方程。
			\[
				\frac{4}{3}\pi r^{3} = \frac{32}{3}\pi.
			\]
			\par\textbf{理由：} 球体积公式 $V = \frac{4}{3}\pi r^{3}$。
		\item[\textbf{第二步（解出 $r$）：}]
			解方程得 $r=2$。
			\[
				r = 2.
			\]
			\par\textbf{理由：} 方程两边除以 $\frac{4}{3}\pi$ 得 $r^{3}=8$，所以 $r=2$。由定理 $r=R$，故底面半径 $R=2$。
		\item[\textbf{第三步（求高）：}]
			利用 $h=2R$ 求高。
			\[
				h = 4.
			\]
			\par\textbf{理由：} 代入 $R=2$ 于 $h=2R$ 得 $h=4$。
	\end{description}
	\textbf{\textcolor{CExample}{关键结论：}}
	\textbf{内切球体积 $\to$ 半径 $r$ $\to$ $R=r$ $\to$ $h=2R$。}

	\medskip
	\textbf{\textcolor{CExample}{例3（综合应用）}}
	\textbf{\textcolor{CExample}{题目：}}
	圆柱存在内切球，其轴截面面积为 $36$，求圆柱的表面积。
	\textbf{\textcolor{CExample}{解题步骤：}}
	\begin{description}[style=nextline, leftmargin=2.8em, labelsep=0.5em, itemsep=0.45em]
		\item[\textbf{第一步（内切球条件）：}]
			设底面半径为 $R$，高为 $h$，由内切球条件得 $h=2R$。
			\[
				h = 2R.
			\]
			\par\textbf{理由：} 圆柱存在内切球的充要条件。
		\item[\textbf{第二步（利用轴截面面积）：}]
			轴截面面积 $h \cdot 2R = 36$，代入 $h=2R$ 得 $4R^{2}=36$，解得 $R=3$。
			\[
				R = 3.
			\]
			\par\textbf{理由：} 由 $h=2R$ 得 $(2R)\cdot(2R)=4R^{2}$。因为轴截面面积 $h\cdot2R=36$，所以 $4R^{2}=36$，故 $R^{2}=9$，$R=3$。
		\item[\textbf{第三步（求表面积）：}]
			圆柱表面积 $S = 2\pi R^{2} + 2\pi R h$。代入 $h=2R$ 得 $S = 2\pi R^{2} + 4\pi R^{2}$，即 $S = 6\pi R^{2}$；代入 $R=3$ 得 $S = 6\pi \cdot 9$，即 $S = 54\pi$。
			\[
				S = 54\pi.
			\]
			\par\textbf{理由：} 由 $S = 2\pi R^{2} + 2\pi R h$，代入 $h=2R$ 得 $S = 2\pi R^{2} + 4\pi R^{2}$，合并得 $S=6\pi R^{2}$。再代入 $R=3$ 得 $S = 6\pi \cdot 9$，即 $S = 54\pi$。
	\end{description}
	\textbf{\textcolor{CExample}{关键结论：}}
	\textbf{存在内切球时，圆柱表面积 $S=6\pi R^{2}$。}
\end{examplebox}
